\documentclass[11pt]{article}

\usepackage[margin=1in]{geometry}
\usepackage{hyperref}
\usepackage[round,authoryear]{natbib}

\title{Literature View:\\
Positioning Remittance Non-Neutrality with Returns}
\author{Draft notes}
\date{\today}

\begin{document}
\maketitle

%=============================================================================
\section{Positioning in the literature}
%=============================================================================

When positioning the paper in the literature, emphasize that existing public
finance work often explains why tax remittance may matter in practice.
\citet{slemrod2008remittance} stress avoidance, evasion, and administrative
costs when the statutory remitter differs from the party that ultimately bears
the tax.  \citet{chetty2009salience} show that consumers may underreact to
taxes that are less salient at the point of purchase, so posted-price versus
register-price remittance can change behavior even when statutory incidence is
held fixed.

This paper identifies a distinct, purely structural source of remittance
non-neutrality that survives fully rational, attentive agents.  Under product
returns, the refund price differs under buyer versus seller remittance: a
seller-remitted tax is refunded at the tax-inclusive checkout price, whereas a
buyer-remitted tax is refunded only at the pre-tax list price.  Because the
keep-or-return decision depends on the refund price, the return margin
dynamically reshapes supply, effective margins, and equilibrium outcomes even
when no agent misperceives the tax and no evasion margin is active.  The
contribution is therefore not another behavioral or enforcement story about
who remits, but a market-structure result: returns break the classical
remittance-neutrality benchmark through the economics of the refund margin
itself.

\bibliographystyle{plainnat}
\bibliography{references}

\end{document}
